\chapter{Problemas Recorrentes e Como Evitá-los}

Este capítulo reúne problemas frequentes em TCCs e orientações para preveni-los. A intenção é ajudar o estudante a revisar o trabalho antes que essas falhas apareçam na banca.

\section{Tema amplo demais}

Um tema amplo torna a pesquisa genérica. O estudante tenta falar de tudo, mas não aprofunda nada. Para corrigir, delimite espaço, tempo, população, objeto, material ou enfoque.

Pergunte: qual parte desse tema eu consigo investigar com rigor no tempo disponível?

\section{Problema que não é pergunta de pesquisa}

Muitos trabalhos apresentam problema como assunto, opinião ou justificativa. O problema precisa orientar investigação. Se ele não pode ser respondido pelo método, ainda não está pronto.

\section{Objetivos desconectados}

Os objetivos específicos devem conduzir ao objetivo geral. Quando cada objetivo aponta para direção diferente, o trabalho perde unidade.

Revise se cada objetivo específico aparece depois na metodologia, nos resultados e na conclusão.

\section{Revisão de literatura descritiva}

Uma revisão frágil resume autor por autor, sem articulação. A revisão deve construir debate. Agrupe autores por ideias, conceitos, convergências e divergências.

\section{Metodologia genérica}

Frases como ``foi realizada pesquisa qualitativa'' ou ``foi feita revisão bibliográfica'' são insuficientes. A metodologia precisa explicar procedimentos concretos.

\section{Resultados sem discussão}

Apresentar dados não é discutir dados. A discussão interpreta resultados, relaciona com objetivos e dialoga com literatura.

\section{Conclusão sem resposta}

A conclusão deve responder ao problema. Se ela apenas repete a introdução, falta fechamento analítico.

\section{Citações em excesso}

Excesso de citações pode apagar a voz do estudante. Use fontes para sustentar o argumento, mas conduza o texto com análise própria.

\section{Referências inconsistentes}

Referências incompletas passam impressão de descuido e dificultam verificação. Faça conferência cruzada: texto para lista final e lista final para texto.

\section{Elementos visuais sem função}

Figuras e tabelas precisam ser necessárias. Se o texto não comenta o elemento, talvez ele não deva estar ali.

\section{Normalização deixada para o fim}

Normalizar apenas na última etapa pode revelar problemas difíceis de corrigir. Ainda assim, a revisão final deve ser feita no PDF, porque quebras de página e listas só aparecem plenamente na versão final.

\section{Falta de revisão após correções}

Alterações feitas depois da banca podem criar erros novos. Sempre revise sumário, paginação, referências e coerência depois de qualquer mudança significativa.

\caixaorientacao{A melhor prevenção é trabalhar por versões. Cada versão deve resolver um tipo de problema: estrutura, conteúdo, método, citações, linguagem e apresentação.}

